\section{Special Abilities}

\subsection{Master Calculator}

Once per turn, for each base with this ability, choose one of the three effects below:

\subsubsection{Trajectory Calculation}

Choose a single detachment within 45~cm of the base with the Master Calculator ability during the target detachment's activation in the Movement Phase. The detachment may use the Master Calculator's line of sight to make its artillery shots. In addition it may reroll the dice it wishes on a Titan or Praetorian hit-location table.

\subsubsection{Structural Weakness}

Choose a building in line of sight of the base with the Master Calculator ability during its activation. Successful save rolls for the building caused by shots from your army may be rerolled.

\subsubsection{Adaptation}

The base gains the Tactical Genius ability for the turn.

\subsection{Versatile}

A base with this ability must choose at the start of the turn (on its activation in the Movement Phase) whether it will have the Mechanic or Medic ability for that turn.

\subsection{Artisan}

A base with this ability acts as a Medic and a Mechanic for Myrmidon, Servitor, and Kataphron bases.

\subsection{Plasma Weapons}

Because these weapons use a large amount of energy, only one Plasma weapon is allowed per titan.

\subsection{Devotion Bell}

All Daemons within 25~cm have a $-2$ penalty to their Instability tests. In addition, no base within 25~cm may receive or use Chaos Reward cards.

\subsection{Ursus Claw}

A weapon with this ability reduces the AF of one target (and only one) in contact by 3.

\subsection{Harpoon Missile}

If the Harpoon Missile hits its target and the target fails its armour save, it remains lodged. While it is lodged, its target suffers a penalty of 3 to AF and a penalty of 1 to To-Hit with all its weapons. In addition to the initial damage, it causes a damage during End-of-Turn Effects while it is lodged. After causing this damage, it dislodges on a 4+ on the first turn it is lodged, 3+ on the second, and 2+ on subsequent turns.

\subsection{Volkite Deflagration}

When a base with this ability shoots, draw a single imaginary line from the base to the aimed point. All bases crossed by that line are hit, whether the centre is hit or not.

These weapons cannot make Intercept Fire. For a structure, half of the bases present in garrison are hit. Due to its nature, it stops at the first structure it encounters as well as the first class 5 or higher obstacle.

\subsection{Reconnaissance Land Speeder}

The Land Speeder is a separate detachment and need not maintain detachment coherency with the titan. If the Titan is destroyed, the Land Speeder is removed from the table. The Land Speeder cannot capture objectives.

\subsection{Crash}

When the Landing Craft is destroyed at Altitude it crashes. It then scatters 2D6~cm from its arrival point, and acts as a Titan falling. If that is the case, the embarked troops lose all their shields. After that, the bases suffer a hit with AP $-2$ for each Wound they have. If they survive they may exit normally from the wreck of the vessel.

\subsection{Vortex Missile}

The Vortex Missile uses a Template (7.5~cm) which, once placed, scatters 3D6~cm. The Vortex Missile does not hit targets at Altitude. It cannot make Intercept Fire. Any base hit by the template (even if the centre is not under the template) suffers the effects of the Vortex Missile. If the target has shields, the vortex may destroy them. The target immediately loses 2D3 Energy Shields and Energy Fields, or 1D3 Deflector Shields and Organic Metal (with no Psychic Save linked to these protections).

Bases that passed their Psychic Save or that still have active shields suffer no damage. All bases are pushed to the edge of the template.

All bases hit suffer 1D3 hits that are Psychic Attacks, and structures are hit by a Damages Buildings (AP $-5$/3) attack.

During End-of-Turn Effects, the vortex disappears on a 5+. If it remains in play, it may be activated in the Combat Phase and then scatters 3D6~cm, hitting everything in its path.

\subsection{Landing Craft}

Landing Craft follow the following arrival rules:

\begin{itemize}
    \item During the Off-Table Arrival phase, the arrival point is determined as with any other off-table arrival.
    \item The vessel arrives in the Movement Phase, and it may suffer Intercept Fire. The vessel is considered to be at Altitude above the arrival point for range and line of sight.
    \item If the vessel survives, it lands and the embarked troops then exit as from a transport. Once on the ground, the vessel can no longer move for the rest of the battle.
\end{itemize}

A Landing Craft may transport 60 points as follows. A Heavy Titan costs 60 points, a Warlord, Psi-Titan Warlord, or Warmaster costs 30 points, a Reaver or Warbringer costs 20 points, a Scout Titan costs 15 points, a Super-heavy costs 10 points, a Knight costs 5 points, a vehicle costs 3 points, a Walker or Cavalry costs 2 points, and an infantry base costs 1 point.

\subsection{Hellbore Emergence}

When the tunneler surfaces, the ground around it is transformed into molten slag. To represent this, place a Template (12~cm) on the tunneler's exit point. All bases under the template are hit on 3+ with AP $-1$. Surviving bases are placed on the edges of the template, as close as possible to their original positions.
